\begin{frame}{P4 -- Vì sao ATPG tuần tự khó?}
\small
\begin{itemize}
  \item DFF lưu trạng thái: $Q(t+1)=D(t)$; $Q(0)$ có thể chưa biết.
  \item Không scan: khó điều khiển $Q$, khó quan sát $D$; cần chuỗi PI
        qua nhiều chu kỳ để khởi tạo, kích hoạt và lan truyền lỗi.
  \item Mạch ví dụ P4: 2 PI, 1 PO, 6 cổng tổ hợp, 1 DFF.
\end{itemize}
\vspace{1mm}
\centering
\resizebox{0.7\textwidth}{!}{\begin{tikzpicture}[>=Latex, node distance=13mm, font=\small]
  \node[draw, rounded corners, minimum width=36mm, minimum height=18mm, align=center] (logic) {Logic tổ hợp\\$N_1,N_2,N_3,D,N_4,Y$};
  \node[draw, rectangle, right=30mm of logic, minimum width=16mm, minimum height=16mm, align=center] (ff) {DFF\\$Q$};
  \draw[->] ([xshift=-18mm,yshift=5mm]logic.west) node[left] {$A$} -- ([yshift=5mm]logic.west);
  \draw[->] ([xshift=-18mm,yshift=-5mm]logic.west) node[left] {$B$} -- ([yshift=-5mm]logic.west);
  \draw[->] ([yshift=5mm]logic.east) -- node[above] {$D$} ([yshift=5mm]ff.west);
  \draw[->] ([yshift=-5mm]logic.east) -- ++(15mm,0) node[right] {$Y$};
  \draw[->] (ff.north) -- ++(0,10mm) -| node[pos=.25,above] {$Q$} (logic.north);
\end{tikzpicture}
}
\end{frame}

\begin{frame}{P4 -- Trải khung và ví dụ $N_3/SA0$}
\small
\begin{itemize}
  \item Thay DFF bằng ràng buộc $Q_{t+1}=D_t$; cùng một lỗi hiện diện
        tại mỗi khung của mạch đã trải.
  \item Dùng hai khung: $(A_0,B_0)=(1,0)$, $(A_1,B_1)=(1,0)$.
\end{itemize}
\begin{center}
\begin{tabular}{@{}c c c c@{}}
\toprule
Khung & $N_3$ tốt/lỗi & $D$ tốt/lỗi & $Y$ tốt/lỗi\\
\midrule
0 & $1/0$ & $0/1$ & $Q_0/Q_0$\\
1 & $1/0$ & $0/1$ & $0/1$\\
\bottomrule
\end{tabular}
\end{center}
\centering Lỗi quan sát tại PO ở khung 1, với cả $Q_0=0$ và $Q_0=1$.
\end{frame}

\begin{frame}{P4 -- Full scan và kết quả kiểm chứng}
\small
\begin{itemize}
  \item Scan FF: shift in $Q$ $\rightarrow$ capture $D$ $\rightarrow$ shift out.
        ATPG tổ hợp xem $Q$ là pseudo-PI, $D$ là pseudo-PO.
  \item Với $N_3/SA0$: $A=1,B=0,Q=0$ cho $D=0/1$.
\end{itemize}
\begin{center}
\begin{tabular}{@{}lcc@{}}
\toprule
Mô hình & Fault universe & Phát hiện\\
\midrule
Full scan, bộ vector vét cạn & 18 stem SA & 18/18\\
Không scan, tối đa 4 khung & 18 stem SA & 18/18 (tối đa 3 khung)\\
\bottomrule
\end{tabular}
\end{center}
\footnotesize Số liệu từ vét cạn và mô phỏng độc lập trên mạch nhỏ;
chưa phải kết quả chạy PODEM của P5/P6.
\end{frame}
